%%%%%%%%%%%%%%%%%%%%%%%%%%%%%%%%%%%%%%%%%%%%%%%%%%%%%%%%%%%%%%%%%%%%%%%%%%%%%%%%
%%
\cleardoubleoddpage%  Make sure to start each chapter on a new odd page
\chapter{Data and Preprocessing}
\label{ch:data}


\section{Dataset}
\label{sec:data:dataset}

The data used in this thesis is a longitudinal cohort of eyes with
Geographic Atrophy, provided by the Medical University of Vienna (MUW).
It was acquired in routine clinical practice at a single clinical site
in Vienna. The sample unit throughout is the individual eye. The cohort
contains 75 eyes, which together contribute 553 visits, where a visit is
one OCT acquisition. Ordering the visits of each eye in time yields 478
visit-to-next-visit intervals ($553 - 75$). The number of visits per eye
ranges from five to thirteen, with a median of seven. The follow-up
span, from the baseline visit to the last recorded visit of an eye,
ranges from 720 to 2160 days, with a median of 1080 days, or roughly
three years.

The acquisition setting recorded in the data is the same for all 553
visits. The vendor is logged as Heidelberg Engineering, the protocol as
a volume scan in ART mode, and the study identifier as ``HRF SPEC''. The
device model is not a recorded field.
% TODO: cite Mai2024 (Mai, Lachinov, Reiter, Riedl, Grechenig, Bogunovic
% & Schmidt-Erfurth, "Deep Learning-Based Prediction of Individual
% Geographic Atrophy Progression from a Single Baseline OCT",
% Ophthalmology Science 4(4):100466, 2024)
Mai et al. (2024), who describe the same MUW cohort, report that it was
acquired on a Spectralis device. The scanner is therefore taken to be a
Heidelberg Spectralis, but this is an inference from that source rather
than a property of the data.
% TODO: confirm the device model (Heidelberg Spectralis) with the MUW
% data provider; it is currently an inference from Mai et al. (2024).

The splits (\S\ref{sec:data:splits}) are defined over a universe of 100
patient identifiers. Of these, 49 have no imaging data. The remaining 51
patients contribute the 75 eyes: 24 patients contribute both eyes and 27
contribute a single eye. Of the 75 eyes, 47 belong to female and 28 to
male patients. Age at baseline ranges from 60.7 to 88.0 years, with a
mean of 76.4 years over the eyes. Ages are decimal values computed from
the date of birth to the baseline visit, not hand-entered integers.

The data does not contain the raw three-dimensional OCT volumes.
Instead, each visit provides a set of registered two-dimensional en-face
derivatives. The first is a binary GA mask on the native OCT en-face
grid, stored as an image with the values 0 and 255. The second is an
array of layer-boundary depth maps of shape (rows, columns, 10), stored
as 16-bit integers in units of axial pixels. Each visit further provides
a GA mask, a scanning laser ophthalmoscopy (SLO) fundus image and a
fundus-autofluorescence image, all in the SLO reference frame, together
with two files of coordinate transforms between the OCT and the SLO
frame. Only the OCT-grid mask and the layer array are required. An eye
is included if at least two of its visits carry both, which guarantees
at least one transition. All 75 eyes meet this condition.

The native en-face grids are not uniform. The 553 visits have 67
distinct grid shapes, with 38 to 66 rows and 961 to 1719 columns. The
shape is constant across the visits of one eye but differs between eyes.
The variation stems from the raw acquisitions, whose B-scan spacing
ranged from 115 to 258\,\textmu{}m. Upstream of this work, the stored
en-face arrays were resampled onto one gold-standard pixel spacing:
0.12118\,mm between B-scans, 0.00568\,mm between A-scans, and
0.003867\,mm per axial pixel. Pixel spacing is therefore treated as
constant across the cohort, while pixel counts, and with them the
physical extent of each native grid, differ between eyes.

% TODO: settle the constant-spacing assumption by obtaining the original
% DICOM spacing fields from the data provider.
This constant spacing is an assumption rather than a measurement.
Checked against the visits' own transform files, it holds to within
about 2\,\%. Settling it requires the original DICOM spacing fields,
which are not part of the data used here. Every area reported in square
millimetres in this thesis inherits this assumption.

Lesion growth is summarised on the square-root area scale.
% TODO: cite Feuer2013 (Feuer, Yehoshua, Gregori, Penha, Chew, Ferris,
% Clemons, Lindblad & Rosenfeld, "Square Root Transformation of
% Geographic Atrophy Area Measurements to Eliminate Dependence of Growth
% Rates on Baseline Lesion Measurements", JAMA Ophthalmology
% 131(1):110-111, 2013)
This transform is standard for GA, because it removes the dependence of
growth rates on baseline lesion size (Feuer et al., 2013). Over the 75
eyes, the median square-root-area growth rate is 0.234\,mm per year. The
interquartile range spans 0.153 to 0.340\,mm per year, and the 90th
percentile is 0.473\,mm per year.
% TODO: add a cohort-level baseline lesion-area distribution table
% comparable to the one reported by Mai et al. (2024).

Figure~\ref{fig:data:example-state} shows one visit of the cohort in the
views used throughout this chapter.

% TODO: produce Figure fig:data:example-state --- three panels for one
% representative visit, all generated from files that exist in the
% repository (the raw B-scans are NOT part of the repository data; a
% B-scan panel would have to be sourced from MUW separately):
% (a) the SLO fundus image with the OCT field of view drawn as a
%     rectangle and the SLO-frame GA mask overlaid;
% (b) channel 0, the binary GA mask on the 49 x 1024 modelling grid,
%     drawn at the physical aspect ratio (about 5.94 x 5.82 mm), so
%     that the ~21:1 pixel anisotropy is visible as tall, thin pixels;
% (c) one or two layer-boundary depth maps (channels 1--10) on the same
%     grid as heatmaps coloured by axial depth, with the GA mask
%     outline from (b) overlaid.
% Any padded region should be marked, since it is visible in about one
% visit in four.
\begin{figure}[t]
  \centering
  \fbox{\parbox[c][6cm][c]{0.9\textwidth}{\centering\itshape Placeholder
    for Figure~\ref{fig:data:example-state}: one visit shown as the SLO
    fundus image with the OCT field of view, the GA mask on the modelling
    grid, and a layer-boundary depth map. See accompanying
    \texttt{\% TODO} for the intended content.}}
  \caption[One visit of the MUW cohort.]{%
    One visit of the MUW cohort.
    \emph{(a)}~SLO fundus image with the OCT field of view marked and the
    GA lesion overlaid.
    \emph{(b)}~The binary GA mask (channel~0) on the $49 \times 1024$
    modelling grid, drawn at its physical aspect ratio of about
    $5.94 \times 5.82$\,mm; each pixel is about 21 times taller than it
    is wide.
    \emph{(c)}~A layer-boundary depth map (one of channels~1--10) on the
    same grid, coloured by axial depth, with the lesion outline from (b)
    overlaid. Panels (b) and (c) show the state tensor as the models
    receive it.}
  \label{fig:data:example-state}
\end{figure}


\section{State Representation}
\label{sec:data:state}

The state of an eye at each visit is an eleven-channel tensor, so
$C = 11$. Channel~0 holds the binary GA mask. Channels~1--10 hold the
ten layer-boundary depth maps, ordered from the most superficial
boundary to the deepest. Channel~0 is both the prediction target and the
clinical deliverable. The layer channels supply structural context and
serve as auxiliary prediction targets.

The tensor is assembled per visit from the two required files. The
stored mask, which has the values 0 and 255, is thresholded above 127.
This yields a binary field with the value 1 inside the lesion and 0
outside it. The 16-bit layer array is cast to floating point. Mask and
layers are then concatenated along the channel axis into a tensor of
shape $11 \times \text{rows} \times \text{columns}$.

The channel order --- the mask first, followed by the ten layers --- is a
fixed invariant of the pipeline. The loss weighting and the layer
encoder both rely on channel~0 being the mask, and the order is checked
by an explicit assertion at load time.

Each layer channel is a depth map. Fixing a pixel of the en-face grid
selects a single A-scan of the underlying volume. The channel value at
that pixel is the axial pixel index at which the boundary surface lies
along that A-scan. Together, the ten channels describe the vertical
layer structure of the retina at every en-face position. Across the
cohort, the raw depth values range from $-14$ to $499$. This slightly
exceeds the nominal 496-pixel depth of the volume, because the upstream
segmentation occasionally returns values just above the top of the
volume or just below its bottom.

On the training split of the canonical fold, the mean depth of the ten
layer channels increases monotonically from 137 to 204 axial pixels. At
0.003867\,mm per axial pixel, this corresponds to a mean depth of about
0.53 to 0.79\,mm below the top of the scan volume. The anatomical names
of the ten boundary surfaces are not recorded in the data as delivered;
the channels carry generic identifiers only.
% TODO: obtain the anatomical names of the ten layer-boundary surfaces
% from the MUW segmentation-pipeline documentation.

Retinal thickness, the distance between the deepest and the shallowest
boundary, is not computed anywhere in this pipeline. It could be
recovered from the layer channels, but it is neither stored nor used as
an input. The reason for keeping all ten layer channels next to the mask
was given in~\S\ref{sec:background:ga-oct:state-rep}: documented
predictors of GA progression are found throughout the vertical retinal
column, not only inside the lesion. The layer geometry is therefore kept
in full rather than reduced to the lesion mask alone. How the channels
are processed by the models is described
in~\S\ref{sec:method:multichannel}.


\section{Spatial Standardisation and its Measured Cost}
\label{sec:data:spatial}

Because the native en-face grids differ in size, every state tensor is
brought to one common grid of 49 rows by 1024 columns. Per axis, a
symmetric centre crop is applied when the native size is at least the
target size, and a symmetric zero-pad is applied when it is smaller.
Both cases are common: 149 of the 553 visits are padded along the rows,
because they have fewer than 49 rows, and 115 are padded along the
columns. Cropping and padding are used instead of resampling for two
reasons. They introduce no interpolation artefacts, and they keep the
stored pixel spacing identical across the cohort by construction. The
use of a $6 \times 6$\,mm analysis window in the literature was reviewed
in~\S\ref{sec:background:ga-oct:muw}.

The resulting window spans $49 \times 0.12118$\,mm by
$1024 \times 0.00568$\,mm, or about $5.938 \times 5.816$\,mm, and holds
50\,176 grid positions. The spacing between B-scans is 21.335 times
larger than the spacing along them. The grid is therefore strongly
anisotropic: moving by one row covers about the same physical distance
as moving by 21 columns. This property matters for every
spatial operator applied to the grid and is taken up again in
Chapter~\ref{ch:method}.

Zero-padding raises a less obvious problem. Physically, a zero means
``no signal'': a mask value of 0 means not atrophic, and a layer depth
of 0 is the top of the volume, above the retina. The models, however, do
not see raw values. Cropping and padding run before the normalisation
statistics are computed (\S\ref{sec:data:normalization}), so the padded
zeros enter the per-channel mean and standard deviation. After
normalisation, the zero fill of the layer channels sits between $-2.35$
and $-2.94$ standard deviations, which is an extreme value. Unless the
pad positions are excluded from the loss, the models are trained to
reproduce this value at full weight. Zero also cannot serve as a
sentinel that marks padding, because real layer values range from $-14$
to $499$ and so include zero.

The crop also has a measured cost for the pathology itself. A census
over all 553 visits shows that the $49 \times 1024$ window cuts off real
lesion area in 27.3\,\% of visits. It removes more than 5\,\% of the
lesion in 6.0\,\% of visits, and 25.7\,\% of the lesion in the worst
case. In addition, 31.1\,\% of the cropped lesions touch the crop
border. Growth across that border is censored, and no missingness flag
records it. The models are therefore trained on censored targets, and
the change-region metrics at the one-year anchor under-measure
progression for roughly one visit in five. The censoring also limits
comparability with literature results computed over the whole field of
view.

The window was nevertheless kept on purpose. Larger windows reduce the
truncation, but they require up to 35\,\% zero-padding and about 1.9
times as many grid positions. They also cannot remove the problem: for
about 19\,\% of lesions, border contact is unavoidable with any window,
because those lesions reach the edge of the scan itself. The
consequences are therefore handled in the evaluation instead
(\S\ref{sec:experiments:protocol}). There, the change-region metric is
split into border-touching and interior lesions, a pad-masked variant of
the metric is reported, and layer errors are computed over non-padded
positions only. In addition, the training loss can optionally exclude
pad positions.


\section{Normalisation}
\label{sec:data:normalization}

After spatial standardisation, each channel is normalised. The
statistics are computed per channel over the training visits of the
current fold only, and are then applied to every visit, including the
validation visits. No validation data enters the statistics. Because
they are computed on the cropped and padded states, the padded zeros
contribute to each channel's mean and standard deviation
(\S\ref{sec:data:spatial}).

The canonical scheme is a z-score,
\[
  x' \;=\; \frac{x - \mu_c}{\sigma_c},
\]
where $\mu_c$ and $\sigma_c$ are the training mean and standard
deviation of channel~$c$. A vanishing standard deviation is replaced by
one to avoid division by zero. A min-max scaling to $[-1, 1]$,
\[
  x' \;=\; 2\,\frac{x - \min_c}{\max_c - \min_c} \;-\; 1,
\]
is kept as an ablation variant.

The binary mask is z-scored in the same way as the layer channels. On
the canonical fold, the background value 0 maps to $-0.602$ and the
foreground value 1 maps to $+1.660$. The binarisation threshold of 0.5
maps to $+0.529$. Two consequences follow. First, the normalised mask
has a hard step at the lesion edge. Second, evaluation maps every
threshold into normalised space rather than mapping the prediction back
to the original scale, and this mapping must be applied consistently.
Table~\ref{tab:data:norm-stats} lists the statistics for the canonical
fold.

\begin{table}[t]
  \centering
  \caption[Per-channel normalisation statistics.]{%
    Per-channel normalisation statistics (mean and standard deviation),
    computed over the training split of the canonical fold after
    cropping and padding. The mask statistics are derived from the two
    mapped values given in the text and rounded to three decimals; the
    layer statistics are still to be extracted.}
  \label{tab:data:norm-stats}
  \begin{tabular}{lrr}
    \toprule
    Channel & Mean & Standard deviation \\
    \midrule
    0 (GA mask)  & 0.266 & 0.442 \\
    1 (layer 0)  & -- & -- \\
    2 (layer 1)  & -- & -- \\
    3 (layer 2)  & -- & -- \\
    4 (layer 3)  & -- & -- \\
    5 (layer 4)  & -- & -- \\
    6 (layer 5)  & -- & -- \\
    7 (layer 6)  & -- & -- \\
    8 (layer 7)  & -- & -- \\
    9 (layer 8)  & -- & -- \\
    10 (layer 9) & -- & -- \\
    \bottomrule
  \end{tabular}
\end{table}
% TODO: extract the per-channel mean and standard deviation of the ten
% layer channels from the precompute metadata (norm_params) of the
% canonical fold and fill Table tab:data:norm-stats. Layer means should
% increase monotonically from about 137 to 204 axial pixels.


\section{Patient-Level Covariates}
\label{sec:data:covariates}

Each window is accompanied by two patient-level covariates, age and sex.
They come from a clinical covariates table with 187 rows, which covers
the wider study cohort and all 75 eyes used here. The table is keyed per
eye, so a patient who contributes two eyes appears in two rows with the
same age and sex. Sex is encoded as 0 for female and 1 for male. A
missing value is imputed with the male fraction of the training split.

Age is z-scored with training statistics. Two encodings are implemented.
In the baseline mode, each eye keeps its baseline age for all of its
windows. In the per-visit mode, which is the canonical setting, each
window carries the age at its own start visit, that is, the baseline age
plus the elapsed years. The z-score statistics for this mode are
computed over all training window-start ages. Since there is one start
age per window, the final visit of each training eye does not enter
them. On the canonical fold this gives 378 ages, with a mean of 77.653
years and a standard deviation of 7.220 years.

In the per-visit mode, the age covariate advances along a rollout. This
leaks no information about the future, because the age at the next step
follows from the current age and the input $\Delta t$
alone.\footnote{The conversion of elapsed days into years uses 365.25
days per year for the age covariate but 365 days per year for $\Delta
t$. The resulting difference is at most 0.07\,\%. It is deliberately
left unchanged, because aligning the two would invalidate every
precomputed dataset.}

The source clinical table also records, per eye, the progression rate,
the pseudodrusen status, the predominant fundus-autofluorescence pattern
and the concentration of hyperreflective foci. None of these variables
is used in this thesis; their possible use is discussed
in~\S\ref{sec:discussion:future}. Whether age and sex contribute to the
predictions is measured in~\S\ref{sec:experiments:ablations}.


\section{Temporal Structure}
\label{sec:data:temporal}

The longitudinal sequences are split into transition windows. For each
eye, the visits are sorted by day in study, and every pair of
consecutive visits forms one window. Window~$i$ consists of the state at
visit~$i$, the state at visit~$i+1$, and the interval $\Delta t$ between
them. Each window holds exactly one transition; this is a fixed
commitment of the framework, not a configuration option. $\Delta t$ is
stored in years, computed as the gap in days divided by 365, so a gap of
180 days gives $\Delta t = 0.4932$. Each window also stores the raw gap
in days, the two covariates and the position of the window within its
eye. Over the cohort this gives 478 windows.

The distribution of the intervals between consecutive visits is given in
Table~\ref{tab:data:intervals}.

\begin{table}[t]
  \centering
  \caption[Distribution of visit intervals.]{%
    Distribution of the intervals between consecutive visits over the
    478 windows of the cohort, with the number of eyes that contain at
    least one interval of each length.}
  \label{tab:data:intervals}
  \begin{tabular}{lrrr}
    \toprule
    Interval & Windows & Share & Eyes \\
    \midrule
    90 days  & 67  & 14.0\,\% & 10 \\
    180 days & 364 & 76.2\,\% & 72 \\
    270 days & 2   & 0.4\,\%  & 2  \\
    360 days & 38  & 7.9\,\%  & 28 \\
    540 days & 6   & 1.3\,\%  & 5  \\
    720 days & 1   & 0.2\,\%  & 1  \\
    \bottomrule
  \end{tabular}
\end{table}

Three properties of this schedule shape the rest of the pipeline. First,
every interval is a multiple of the 90-day acquisition grid, and the
most common interval is 180 days. Second, the short 90-day intervals are
concentrated in a small subgroup: all 67 of them belong to 10 eyes from
7 patients. A training scheme restricted to intervals shorter than the
modal 180 days would therefore learn from this densely imaged subgroup
alone.

Third, the schedule limits how far an eye can be rolled forward during
training, where several visits are chained within a fixed time budget
(\S\ref{sec:method:training}). With a budget of 90 days, none of the 75
eyes can chain two visits. With 360 days, 74 of the 75 eyes reach at
least one autoregressive step. With 540 days, 68 eyes reach at least two
steps, and with 720 days, 73 eyes do.

The same schedule also shapes the evaluation. The one-year anchor has no
upper bound: the metric is taken at the first rollout step whose
cumulative time reaches about 347 days. For 68 of the 75 eyes, this is
exactly 360 days. The remaining 7 eyes (9.3\,\%) are scored at 450 to
720 days, where the true change region is 1.3 to 1.7 times larger.
% TODO: cite Mai2024 (see the marker in Section 3.1).
This is accepted on purpose, for comparability with Mai et al. (2024).
The evaluation protocol is defined in~\S\ref{sec:experiments:protocol}.


\section{Splits and Cross-Validation}
\label{sec:data:splits}

The data is partitioned by five split files that define a patient-level
5-fold cross-validation over the universe of 100 patient identifiers. In
each fold, 80 patients are assigned to training and 20 to validation.
The five validation sets do not overlap, and together they cover the
whole universe. The split identifiers carry no eye suffix, so both eyes
of a patient always fall into the same partition. This prevents leakage
between the eyes of one patient, which matters because GA is frequently
bilateral with strongly correlated eyes
(\S\ref{sec:background:ga-oct:muw}).

Only 51 of the 100 patients have scans. The per-fold eye counts are
therefore uneven, although the patient split is balanced.
Table~\ref{tab:data:splits} lists them. On fold~2, the split yields 378
training windows and 100 validation windows.
% TODO: record the training and validation window counts for folds 0,
% 1, 3 and 4.

\begin{table}[t]
  \centering
  \caption[Eyes per cross-validation fold.]{%
    Number of training and validation eyes in each of the five
    cross-validation folds. Each fold assigns 80 of the 100 patient
    identifiers to training and 20 to validation; the eye counts are
    uneven because 49 of the identifiers have no scans.}
  \label{tab:data:splits}
  \begin{tabular}{lrr}
    \toprule
    Fold & Training eyes & Validation eyes \\
    \midrule
    0 & 55 & 20 \\
    1 & 63 & 12 \\
    2 & 59 & 16 \\
    3 & 60 & 15 \\
    4 & 63 & 12 \\
    \bottomrule
  \end{tabular}
\end{table}

The test split is structurally empty. An eye whose patient appeared in
neither the training nor the validation list would be assigned to a test
partition. However, every patient with scans lies inside the
100-identifier universe, and every identifier is assigned to training or
validation in every fold. No test partition is therefore ever produced.
The protocol is 5-fold train/validation cross-validation, and every
result reported in this thesis is a validation result.
% TODO: ensure that no chapter describes the cross-validation results as
% test-set performance; the test split is structurally empty.

The folds also differ in baseline lesion size. The median baseline
lesion area is about 4.1\,mm$^2$ on fold~4, against 6.0 to 7.9\,mm$^2$
on the other four folds. This heterogeneity between folds is revisited
in the evaluation protocol (\S\ref{sec:experiments:protocol}).


%%
%%%%%%%%%%%%%%%%%%%%%%%%%%%%%%%%%%%%%%%%%%%%%%%%%%%%%%%%%%%%%%%%%%%%%%%%%%%%%%%%
